
\documentclass[12pt]{article}
\usepackage[utf8]{inputenc}
\usepackage{graphicx}
\usepackage{subcaption}
\usepackage{amsmath}
\usepackage{geometry}
\usepackage{setspace}
\usepackage[hidelinks]{hyperref}
\usepackage{booktabs}
\usepackage{siunitx}
\usepackage{float}
\usepackage{longtable}

\geometry{margin=1in}
\setstretch{1.5}

\usepackage[backend=biber,style=ieee,sorting=nyt]{biblatex}
\addbibresource{references.bib}

\title{Reference-Based Pose Evaluation for Automated Martial Arts Technique Assessment}
\author{Jimmy Huynh}
\date{March 2026}

\begin{document}

\pagenumbering{gobble}
\maketitle
\thispagestyle{empty}
\clearpage
\pagenumbering{arabic}

\begin{abstract}
This thesis presents the design and evaluation of a reference-based system for automated assessment of martial-arts techniques from video. The system uses pose estimation to extract skeletal keypoints and compares user-performed movements against technique-specific reference sequences to produce a continuous score and structured feedback.

The project investigates whether such a pose-based trainer can distinguish between stronger and weaker executions and move toward agreement with expert judgement. The implemented system includes reference-pose generation, clip-level scoring, threshold-based calibration support, and analysis utilities for comparing model outputs across techniques and users.

The evaluation shows that reference-based pose comparison is a viable foundation for automated martial-arts training support and that the system can produce interpretable clip-level assessments. At the same time, the results indicate that stronger validation requires broader participant-level testing, labeled calibration data, and more systematic comparison with expert assessment.

Overall, the thesis demonstrates both the practical potential and the current limitations of using pose-based reference matching for technique evaluation in martial-arts training.
\end{abstract}

\newpage
\section{Background and Literature Review}
\subsection{Human Action Recognition and Pose Estimation}
Human Action Recognition (HAR) seeks to identify movements or activities from image or video data. In sports contexts, HAR is especially useful because it can transform raw video into structured information about technique, and performance. However, sports video is difficult to analyse due to occlusion, variable viewpoints, large appearance differences, and rapid motion. These difficulties are even more pronounced in combat sports, where a technique may last only a fraction of a second and body parts are often partially hidden \cite{wu2024sthka}.

Pose estimation addresses part of this problem by reducing video to a set of semantically meaningful keypoints. Instead of relying on raw pixels, a pose-based method represents the athlete as a skeleton, making it easier to compare body structure over time and across performers. Prior work indicates that pose estimation is especially relevant for martial arts and karate because technique quality is closely linked to joint alignment, body mechanics, and coordinated motion \cite{echeverria2021pose}. This makes pose-based analysis a suitable basis for a coaching-oriented system.

\subsection{Combat Sports Analytics}
Recent research has explored the use of computer vision for combat sports scoring, athlete tracking, and action recognition. Quinn and Corcoran argue that real-time object detection and pose estimation can be combined to generate meaningful performance statistics in combat sports and potentially reduce the effects of human error in scoring \cite{quinn2024automation}. Similarly, work on optimized YOLO pipelines for combat sports highlights the importance of careful data collection, data augmentation, modular system design, and domain-specific optimization when working with athletic video data \cite{quinn2025yolo}.

Beyond officiating and statistics, related work also points toward coaching support. Pose-estimation systems are attractive in athletic settings because they can turn a movement into joint trajectories, timing information, and body-angle measurements that are easier to compare against a target pattern than raw video alone \cite{echeverria2021pose,wu2024sthka}. In practice, this means such systems can be used not only to recognise that an athlete performed a jab or kick, but also to indicate where the execution differs from a technical reference. That is a more useful coaching scenario than generic classification, because athletes need actionable feedback on extension, guard position, timing, and recovery rather than only a label.

The literature also suggests an important usability requirement for sports-training systems: the feedback must be understandable and tied to a training context. A large generic model may recognise many motions, but it can be difficult for a coach to know what technical standard the model is implicitly applying. For this reason, reference-based approaches are particularly attractive in athlete support settings. They allow a coach or domain expert to define the target movement explicitly through curated exemplars, viewpoints, and technique categories. This makes the resulting assessment more transparent, easier to adapt to a specific gym or style, and easier to explain to athletes during training.

At the same time, the literature makes clear that generic video classification alone is not enough. A technically useful system must cope with viewpoint changes, noisy data, short action windows, and the absence of large labelled datasets specific to combat sports \cite{wu2024sthka,quinn2025yolo}. These issues strongly influenced the design decisions in the project.

In recent years, several real-world systems and commercial applications have emerged to support athletes using computer vision and pose estimation. For example, the BlazePose model, developed by Google, is integrated into fitness apps and wearables to provide real-time feedback on exercise form and posture, enabling users to correct their movements during workouts \cite{wu2024sthka}. Commercial platforms such as Coach's Eye and Dartfish offer video analysis tools that allow coaches and athletes to annotate, compare, and review technique frame-by-frame, although these typically rely on manual input rather than automated pose extraction. More advanced systems, like the MoveNet and OpenPose frameworks, have been incorporated into prototype coaching applications that automatically detect keypoints and provide feedback on joint angles, symmetry, and movement efficiency.

Research prototypes have also demonstrated the potential of automated feedback in sports training. Echeverria and Santos \cite{echeverria2021pose} developed a system for karate that uses pose estimation to model psychomotor performance, providing quantitative metrics for technique assessment. Similarly, Quinn and Corcoran \cite{quinn2024automation} describe a real-time combat sports analysis platform that combines object detection and pose estimation to generate performance statistics and reduce human error in judging. These examples illustrate the growing adoption of computer vision-based systems for athlete support, ranging from commercial video review tools to fully automated feedback platforms in both research and practice.

\section{Experimental Setup}
The uploaded Git Repository supports two related evaluation modes. The first is \textit{reference collection and calibration}, where canonical technique examples are captured and validated. The second is \textit{user evaluation}, where an athlete performs a target technique and the system compares that pose sequence against the calibrated reference bank.

Reference examples are stored as NumPy arrays of keypoint sequences in COCO-17 format \cite{projectActionRecognition}. The project documentation specifies a multi-angle folder structure for techniques, while the progress document shows that the intended capture plan consists of 52 total targets. The batch collection script aims by default for four examples per angle, subject to sufficient source diversity \cite{projectRunBatch}. This means the intended experiment design is systematic rather than ad hoc.

For evaluation, \texttt{calibrate\_label\_thresholds.py} requires validation examples for each technique split into \texttt{correct} and \texttt{incorrect} folders \cite{projectCalibrate}. Every validation clip is scored using the same best-reference matching routine used by the trainer, and the threshold is selected automatically from the observed score distribution. This is important because it means the acceptance rule is derived from validation evidence rather than hand tuning.

For an \textit{amateur versus professional} experiment, both groups can perform the same target techniques under comparable recording conditions, after which the system score, threshold decision, and score variance can be compared across groups. For an \textit{expert versus system} experiment, professional practitioner provide an oral score on a 1--100 scale for each attempt, and that score can then be compared against a normalized version of the system's internal score. The code supports the core metrics needed for these experiments.

\section{Challenges and Solutions}
\subsection{Challenge 1: Building a Reference Dataset}
One of the most difficult parts of the project was not model inference itself, but the creation of a reliable reference library. Unlike benchmark-driven projects that begin with a large labelled dataset, this system depends on curated canonical examples of techniques. This aligns with the literature, which identifies data collection, annotation, and dataset quality as major constraints in combat sports computer vision \cite{quinn2024automation,quinn2025yolo}.

The codebase addresses this challenge through a semi-automated workflow. 

First, \texttt{scrape\_jab\_candidates.py} collects candidate YouTube videos for a target technique using the YouTube Data API and stores them with metadata such as title, duration, and URL \cite{projectScrape}. Second, \texttt{filter\_candidates.py} applies hand-designed keep and reject keyword patterns to filter obvious false positives and non-technique content \cite{projectFilter}. Third, \texttt{generate\_reference\_capture\_commands.py} converts reviewed CSV rows into executable PowerShell commands for reference capture \cite{projectGenerate}. Fourth, \texttt{run\_reference\_collection\_batch.py} executes these capture jobs in batch mode with checks for source diversity, controls on examples per angle, CPU throttling, and cooldown periods \cite{projectRunBatch}. This workflow directly addresses the challenge of collecting high-quality references from messy web video sources.
The most consistent type of data has been extracting self made videos or reviewed youtube videos, as it can be recorded in controlled environments to reduce noise.

\subsection{Challenge 2: Viewpoint Sensitivity}
A technique in 2D pose space can vary substantially with camera placement. A jab observed from the side may appear very different from a jab observed from the front, even when biomechanically correct. Prior literature identifies viewpoint variation as a core difficulty in martial arts and action recognition more broadly \cite{wu2024sthka,echeverria2021pose}.

During development it has been discovered that a side view, the Orthodox (left foot/hand forward, right-handed) fighting stance yields the best results. Technical difficulties can be seen using Soutpaw (right foot/hand forward, left-handed) with right side view and vice-versa, as this causes hidden joints from specific angles. This causes the normalization to misalign center position and the sizing of the pose figure.

To combat this issue, a multi-angle reference design has been implemented. The project makes use of a technique-first folder layout in which each technique can store separate angle files such as \texttt{front}, \texttt{left45}, \texttt{right45} \texttt{left}, \texttt{right} and \texttt{behind }\cite{projectReadme}. During evaluation, the code does not compare an athlete against a single reference template. Instead, it evaluates all available angle references for the target technique and selects the one with the best score \cite{projectActionRecognition}. This substantially improves robustness to viewpoint mismatch.

\subsection{Challenge 3: Temporal Segmentation of Techniques}
A combat technique is not defined by one frame. A punch, for example, includes initiation, extension, peak action, and recovery. If a saved reference window begins too early or ends too late, the resulting sequence contains irrelevant motion and becomes a weak template. This problem is especially important in combat sports because many actions are brief and explosive, but during technical training, often slow movement and prolonged duration.

The project addresses this challenge inside \texttt{action\_recognition.py} through explicit sequence extraction logic. The code supports event-centered extraction around peak wrist extension and a stance-cycle mode based on start, peak, and end thresholds \cite{projectActionRecognition}. It also uses wrist motion energy and wrist return closure as gates before saving a reference, so that only sequences with meaningful motion and plausible return behaviour are accepted \cite{projectActionRecognition}. This is a practical solution to the problem of noisy or poorly segmented reference windows.
Combining this solution with Dynamic Time Warping creates a program, that work for both explosive and slow actions.

\subsection{Challenge 4: Translating Similarity into Correctness}
A raw score is informative, but it does not by itself define what counts as a correct technique. A threshold that is appropriate for one action may be too strict or too lenient for another. This problem becomes more severe when different techniques vary in visibility, temporal stability, and sensitivity to joint noise.

The solution implemented in the project is \texttt{calibrate\_label\_thresholds.py}. This script assumes a validation layout of \texttt{datasets/validation/<technique>/correct/*.npy} and \texttt{datasets/validation/<technique>/incorrect/*.npy}. For each technique, it scores validation sequences against the reference bank, searches across candidate thresholds, and chooses the threshold that maximizes F1 score and then balanced accuracy \cite{projectCalibrate}.

\subsection{Challenge 5: Runtime Constraints and Detection Gaps}
The runtime logs show that pose-based reference collection is feasible on CPU hardware, but not trivial. This project has been run using the following specifications:

\begin{itemize}
    \item Processor	AMD Ryzen 9 5900X 12-Core Processor             (3.70 GHz)
    \item Installed RAM	16,0 GB
    \item System type	64-bit operating system, x64-based processor
    \item NVIDIA GeForce GTX 1660 SUPER
    \item NVIDIA Cuda ON
\end{itemize}

In front-kick reference expansion runs, the code processed frames at approximately 10--16 frames per second in many segments, and the logs also contain repeated \textit{no detections} frames \cite{projectFrontKickLog,projectFrontKickStrictLog}. The seed batch run log similarly illustrates the non-negligible cost of continuous pose inference \cite{projectSeedLog}. This is consistent with the literature, which argues that real-time combat analysis is achievable only with careful management of resources and data quality \cite{quinn2024automation,quinn2025yolo}.

The implementation mitigates this through several engineering controls. The scripts expose configurable frame skipping, CPU thread limits, cooldown periods between batch jobs, preflight source checks, headless operation, and a \texttt{fast-mode} preset that disables expensive visual components \cite{projectActionRecognition,projectRunBatch}. Furthermore, CUDA enabled massive parallelism, allowing tasks to run much faster than on traditional CPUs. This has accelerated video processing significantly. 
These measures are not incidental; they are part of the project's answer to practical deployment constraints.

\section{Results}
This section presents the main evaluation results for the implemented system across supported techniques and user groups.

\subsection{Technique Score Summary}

Figure~\ref{fig:score_by_technique} shows the distribution of scores for all supported techniques, aggregated across all evaluated clips. This summary visualization allows for direct comparison of system performance and threshold calibration across techniques, highlighting which actions are more consistently performed above or below the acceptance threshold.

\subsection{Score Distributions by Technique and User}

\begin{figure}[H]
    \centering
    % Replace 'score_distributions_by_technique.png' with your actual plot file
    \includegraphics[width=0.95\textwidth]{score_by_technique_by_technique.png}
    \caption{Score distributions by technique and user group. Each boxplot shows the spread of scores for a given technique, grouped by user and labeled as (Amateur) or (Professional).}
    \label{fig:score_by_technique}
\end{figure}

Figure~\ref{fig:score_by_technique} presents the distribution of system-generated scores for each technique, separated by user and group. The boxplots reveal several key findings:

\begin{itemize}
    \item \textbf{Discriminative Power:} The system distinguishes between amateur and professional users, with professionals generally achieving higher and more consistent scores, especially for techniques like axe kick and front kick.
    \item \textbf{Technique Variability:} Some techniques, such as jab, show a wider spread of scores, indicating greater variability in performance or higher sensitivity in the scoring algorithm.
    \item \textbf{Outliers and Consistency:} The presence of outliers suggests occasional detection errors or natural performance variation. Professionals tend to have fewer outliers and narrower score ranges.
\end{itemize}

These results demonstrate that the program’s scoring system is sensitive to both user skill and technique, aligning with the intended design of providing meaningful, technique-specific feedback.

\section{Analysis}
The results demonstrate that the reference-based pose evaluation system is capable of distinguishing between stronger and weaker executions of martial arts techniques. Professionals consistently achieve higher scores and exhibit less variance, indicating that the system's scoring aligns with expected skill differences. The use of multi-angle references and temporal alignment methods (such as Dynamic Time Warping) improves robustness to viewpoint and speed variations, while the calibration process ensures that thresholds are technique-specific and evidence-based.

However, the presence of outliers and occasional misclassifications highlights the ongoing challenges of pose estimation accuracy, segmentation, and the need for high-quality reference data. The system's ability to provide interpretable feedback and visual overlays is a practical strength, supporting both automated assessment and actionable coaching.

\section{Discussion}
The findings suggest that reference-based pose comparison is a viable foundation for automated martial-arts training support. The modular pipeline, explicit reference library, and calibration routines make the system transparent and adaptable to new techniques or training contexts. The approach bridges the gap between generic action recognition and domain-specific coaching needs by providing interpretable, technique-aligned feedback.

Nevertheless, the system's performance is limited by the quality and diversity of reference data, the accuracy of pose estimation under challenging conditions, and the inherent variability in human movement. Broader validation, including more participants and expert-labeled calibration data, is needed to further assess reliability and generalizability. The engineering solutions for runtime constraints and viewpoint sensitivity are effective but could be further optimized for real-time deployment.

\section{Conclusion}
This thesis demonstrates both the practical potential and the current limitations of using pose-based reference matching for technique evaluation in martial-arts training. The implemented system can produce interpretable clip-level assessments and distinguish between different skill levels, moving toward agreement with expert judgement. At the same time, the results indicate that stronger validation requires broader participant-level testing, labeled calibration data, and more systematic comparison with expert assessment.

\section{Future Work}
Future work should focus on expanding the participant pool for evaluation, improving the diversity and quality of reference data, and integrating more systematic expert comparison. Enhancements to pose estimation accuracy, real-time feedback mechanisms, and user interface design could further increase the system's utility in practical training settings. Exploring the integration of additional sensor data or multimodal analysis may also provide richer feedback and more robust assessment capabilities.
Assessing martial-arts technique is traditionally a human-centred task. Coaches and judges evaluate movements based on criteria such as timing, balance, control, alignment, and impact, but these assessments can vary across observers and are difficult to apply consistently in video-based review. Recent work in computer vision and pose estimation suggests that athletic movement can be analysed from video using vision-based and skeletal representations~\cite{echeverria2021karate,quinn2024automation,wu2024blazepose}.

This thesis investigates a reference-based approach to automated martial-arts technique assessment. Instead of treating the problem as action classification alone, the project develops a pose-driven trainer that compares a performed movement against technique-specific reference sequences. The implemented system combines person tracking, pose estimation, temporal pose-sequence processing, multi-angle reference matching, and threshold-based scoring~\cite{actionRecognition}.

The central aim is to examine whether such a system can distinguish between stronger and weaker executions of a technique and move toward agreement with expert judgement. To support this, the project implements the technical pipeline required for clip-level evaluation, including reference-pose generation, structured run logging, and calibration support~\cite{calibrateThresholds,runReferenceCollectionBatch}.

The remainder of the report is structured as follows. First, relevant background and related work are reviewed. Next, the system design and implementation are presented, followed by the main development challenges and adopted solutions. The report then presents the evaluation results, discusses their implications and limitations, and concludes with directions for future work.

\section{Background and Literature}
\subsection{Human Action Recognition and Pose Estimation}
Human Action Recognition (HAR) seeks to identify movements or activities from image or video data. In sports contexts, HAR is especially useful because it can transform raw video into structured information about technique, and performance. However, sports video is difficult to analyse due to occlusion, variable viewpoints, large appearance differences, and rapid motion. These difficulties are even more pronounced in combat sports, where a technique may last only a fraction of a second and body parts are often partially hidden \cite{wu2024sthka}.

Pose estimation addresses part of this problem by reducing video to a set of semantically meaningful keypoints. Instead of relying on raw pixels, a pose-based method represents the athlete as a skeleton, making it easier to compare body structure over time and across performers. Prior work indicates that pose estimation is especially relevant for martial arts and karate because technique quality is closely linked to joint alignment, body mechanics, and coordinated motion \cite{echeverria2021pose}. This makes pose-based analysis a suitable basis for a coaching-oriented system.

\subsection{Combat Sports Analytics}
Recent research has explored the use of computer vision for combat sports scoring, athlete tracking, and action recognition. Quinn and Corcoran argue that real-time object detection and pose estimation can be combined to generate meaningful performance statistics in combat sports and potentially reduce the effects of human error in scoring \cite{quinn2024automation}. Similarly, work on optimized YOLO pipelines for combat sports highlights the importance of careful data collection, data augmentation, modular system design, and domain-specific optimization when working with athletic video data \cite{quinn2025yolo}.

Beyond officiating and statistics, related work also points toward coaching support. Pose-estimation systems are attractive in athletic settings because they can turn a movement into joint trajectories, timing information, and body-angle measurements that are easier to compare against a target pattern than raw video alone \cite{echeverria2021pose,wu2024sthka}. In practice, this means such systems can be used not only to recognise that an athlete performed a jab or kick, but also to indicate where the execution differs from a technical reference. That is a more useful coaching scenario than generic classification, because athletes need actionable feedback on extension, guard position, timing, and recovery rather than only a label.

The literature also suggests an important usability requirement for sports-training systems: the feedback must be understandable and tied to a training context. A large generic model may recognise many motions, but it can be difficult for a coach to know what technical standard the model is implicitly applying. For this reason, reference-based approaches are particularly attractive in athlete support settings. They allow a coach or domain expert to define the target movement explicitly through curated exemplars, viewpoints, and technique categories. This makes the resulting assessment more transparent, easier to adapt to a specific gym or style, and easier to explain to athletes during training.

At the same time, the literature makes clear that generic video classification alone is not enough. A technically useful system must cope with viewpoint changes, noisy data, short action windows, and the absence of large labelled datasets specific to combat sports \cite{wu2024sthka,quinn2025yolo}. These issues strongly influenced the design decisions in the project.

In recent years, several real-world systems and commercial applications have emerged to support athletes using computer vision and pose estimation. For example, the BlazePose model, developed by Google, is integrated into fitness apps and wearables to provide real-time feedback on exercise form and posture, enabling users to correct their movements during workouts \cite{wu2024sthka}. Commercial platforms such as Coach's Eye and Dartfish offer video analysis tools that allow coaches and athletes to annotate, compare, and review technique frame-by-frame, although these typically rely on manual input rather than automated pose extraction. More advanced systems, like the MoveNet and OpenPose frameworks, have been incorporated into prototype coaching applications that automatically detect keypoints and provide feedback on joint angles, symmetry, and movement efficiency.

Research prototypes have also demonstrated the potential of automated feedback in sports training. Echeverria and Santos \cite{echeverria2021pose} developed a system for karate that uses pose estimation to model psychomotor performance, providing quantitative metrics for technique assessment. Similarly, Quinn and Corcoran \cite{quinn2024automation} describe a real-time combat sports analysis platform that combines object detection and pose estimation to generate performance statistics and reduce human error in judging. These examples illustrate the growing adoption of computer vision-based systems for athlete support, ranging from commercial video review tools to fully automated feedback platforms in both research and practice.

\section{Experimental Setup}
The implemented project is not a simple end-to-end classifier that takes a video and outputs a label. Instead, it is a modular, pose-driven pipeline built around person tracking, pose-sequence extraction, reference matching, and threshold-based scoring. The central implementation is contained in \texttt{action\_recognition.py}, which uses an Ultralytics YOLO pose model to detect and track people, extract keypoints in COCO-17 format, and maintain temporal histories of those keypoints for each tracked identity \cite{projectActionRecognition}. The code defines support for the labels \textit{fighting stance, jab, cross, hook, uppercut, front kick, roundhouse kick, side kick, back kick, spinning back kick, knee strike, elbow strike,} and \textit{axe kick} \cite{projectActionRecognition}. These labels are not fixed, since the system permits arbitrary movement sequences to be defined and associated with user-specified labels.


The pipeline begins with a video source, which may be a webcam, local file or a YouTube URL. For each frame, the YOLO pose model performs person detection and pose estimation. The extracted poses are associated with persistent track identifiers, allowing the system to follow one athlete over time rather than analysing isolated frames only \cite{projectActionRecognition}. This temporal continuity is important because a technique such as a jab or kick cannot be judged from a single still image.

The overall architecture is best understood as two connected phases rather than one monolithic process. The first phase is a \emph{setup or training phase}, in which the trainer curates the reference bank, captures or selects canonical examples, organises them by technique and angle, and, where data is available, calibrates thresholds. The second phase is the \emph{live analysis phase}, in which a new user performance is recorded, converted into a pose sequence, matched against the curated bank, and scored. Separating these phases is important conceptually: the system does not learn a hidden standard from a large black-box dataset during use, but instead evaluates against an explicit reference library that can be inspected and extended.

This also means the system is not locked to only the current techniques. The present implementation already includes a predefined set of labels and heuristics for common punches and kicks \cite{projectActionRecognition}, but the underlying design is extendable. New techniques can be added by creating additional reference sequences, organizing them under the same technique-first folder structure, and, where needed, defining technique-specific angle or threshold settings. In that sense, the practical limit is not the pipeline itself, but the availability and quality of curated references for the new motion.

\begin{figure}[H]
\centering
\fbox{\parbox{0.94\textwidth}{
\centering
\textbf{Setup / training phase}\\[2mm]
Select target technique $\rightarrow$ collect or record reference clips $\rightarrow$ extract pose sequences $\rightarrow$ organize by technique and angle $\rightarrow$ optionally calibrate thresholds\\[4mm]
\textbf{Live analysis phase}\\[2mm]
Capture user video $\rightarrow$ detect and track pose $\rightarrow$ extract temporal window $\rightarrow$ compare against reference bank $\rightarrow$ choose best-matching variation $\rightarrow$ score, classify, and return feedback
}}
\caption{High-level separation between the offline setup phase and the live analysis phase.}
\label{fig:system_pipeline}
\end{figure}

The implementation also supports optional video classification. The code includes interfaces for pretrained TorchVision video models such as S3D, R3D-18, Swin3D, and MViT, as well as a zero-shot Hugging Face X-CLIP classifier \cite{projectActionRecognition}. However, the project architecture clearly prioritises pose-based scoring. In many of the reference-collection workflows, the video classifier is explicitly disabled so that the system focuses on clean pose extraction and technique comparison instead of generic action labels \cite{projectFrontKickLog,projectSeedLog}.

A central design element is the reference-pose library. Canonical techniques are stored as \texttt{.npy} files containing pose sequences. The loader supports both a legacy flat structure and a multi-angle layout of the form \texttt{reference\_poses/<technique>/<angle>.npy} \cite{projectActionRecognition,projectReadme}. This is an important methodological choice because the same technique can appear substantially different across camera angles. Rather than forcing every user attempt to match a single canonical view or technique variation, the system selects the best-matching angle and movement pattern from the reference bank \cite{projectActionRecognition}.

The specific pose variation is therefore not chosen manually at runtime by the user. Instead, for the requested technique, the system evaluates the available reference variants and selects the one that yields the best overall match score. In practical terms, this means that a jab performed from one viewpoint can be compared against side-view jab references, while the same technique from a different viewpoint can be matched against a more suitable frontal or 45-degree reference if such a variant exists in the bank. The quality of this selection depends directly on how representative the curated reference variations are.

The scoring process begins by normalising each pose frame. The body is centred using the hip midpoint when available, and the coordinates are scaled by shoulder distance or a fallback body-based scale \cite{projectActionRecognition}. This reduces the influence of camera distance, translation, and body size. The system then compares the user sequence against the reference sequence using several pose-based measures. These include Dynamic Time Warping for temporal alignment, cosine similarity of flattened pose trajectories, and technique-specific joint angle error using the joints most relevant to punches or kicks \cite{projectActionRecognition}. For striking actions, the code also computes wrist motion energy and wrist return closure, two handcrafted measures intended to capture whether a movement contains both meaningful extension and a realistic return phase \cite{projectActionRecognition}.

The trainer output is therefore richer than a class label. For each evaluated attempt, the system can return a best-matching angle, a numerical score, a score threshold, a correctness decision, and generated feedback \cite{projectActionRecognition}. If the attempt falls below the threshold, the code constructs a ghost-pose overlay from the reference sequence, allowing visual comparison between the observed movement and the expected technique \cite{projectActionRecognition}. This is particularly relevant in a training setting because it connects the numerical evaluation to a concrete visual correction. The code also includes a pathway for exporting oral scores and aligning them with program outputs, thereby supporting later comparison between expert judgement and system-generated scores, \cite{projectEvaluateOralScores}.

The implementation also includes structured storage of experimental artifacts. It can create dedicated run folders containing configuration metadata, metric rows, track summaries, and overlays \cite{projectActionRecognition}. From a research perspective, this is valuable because it improves reproducibility, makes experiments easier to compare, and supports later analysis of failures and edge cases.

\section{Challenges and Solutions}
\subsection{Challenge 1: Building a Reference Dataset}
One of the most difficult parts of the project was not model inference itself, but the creation of a reliable reference library. Unlike benchmark-driven projects that begin with a large labelled dataset, this system depends on curated canonical examples of techniques. This aligns with the literature, which identifies data collection, annotation, and dataset quality as major constraints in combat sports computer vision \cite{quinn2024automation,quinn2025yolo}.

The codebase addresses this challenge through a semi-automated workflow. 

First, \texttt{scrape\_jab\_candidates.py} collects candidate YouTube videos for a target technique using the YouTube Data API and stores them with metadata such as title, duration, and URL \cite{projectScrape}. Second, \texttt{filter\_candidates.py} applies hand-designed keep and reject keyword patterns to filter obvious false positives and non-technique content \cite{projectFilter}. Third, \texttt{generate\_reference\_capture\_commands.py} converts reviewed CSV rows into executable PowerShell commands for reference capture \cite{projectGenerate}. Fourth, \texttt{run\_reference\_collection\_batch.py} executes these capture jobs in batch mode with checks for source diversity, controls on examples per angle, CPU throttling, and cooldown periods \cite{projectRunBatch}. This workflow directly addresses the challenge of collecting high-quality references from messy web video sources.
The most consistent type of data has been extracting self made videos or reviewed youtube videos, as it can be recorded in controlled environments to reduce noise.

\subsection{Challenge 2: Viewpoint Sensitivity}
A technique in 2D pose space can vary substantially with camera placement. A jab observed from the side may appear very different from a jab observed from the front, even when biomechanically correct. Prior literature identifies viewpoint variation as a core difficulty in martial arts and action recognition more broadly \cite{wu2024sthka,echeverria2021pose}.

During development it has been discovered that a side view, the Orthodox (left foot/hand forward, right-handed) fighting stance yields the best results. Technical difficulties can be seen using Soutpaw (right foot/hand forward, left-handed) with right side view and vice-versa, as this causes hidden joints from specific angles. This causes the normalization to misalign center position and the sizing of the pose figure.

To combat this issue, a multi-angle reference design has been implemented. The project makes use of a technique-first folder layout in which each technique can store separate angle files such as \texttt{front}, \texttt{left45}, \texttt{right45} \texttt{left}, \texttt{right} and \texttt{behind }\cite{projectReadme}. During evaluation, the code does not compare an athlete against a single reference template. Instead, it evaluates all available angle references for the target technique and selects the one with the best score \cite{projectActionRecognition}. This substantially improves robustness to viewpoint mismatch.

\subsection{Challenge 3: Temporal Segmentation of Techniques}
A combat technique is not defined by one frame. A punch, for example, includes initiation, extension, peak action, and recovery. If a saved reference window begins too early or ends too late, the resulting sequence contains irrelevant motion and becomes a weak template. This problem is especially important in combat sports because many actions are brief and explosive, but during technical training, often slow movement and prolonged duration.

The project addresses this challenge inside \texttt{action\_recognition.py} through explicit sequence extraction logic. The code supports event-centered extraction around peak wrist extension and a stance-cycle mode based on start, peak, and end thresholds \cite{projectActionRecognition}. It also uses wrist motion energy and wrist return closure as gates before saving a reference, so that only sequences with meaningful motion and plausible return behaviour are accepted \cite{projectActionRecognition}. This is a practical solution to the problem of noisy or poorly segmented reference windows.
Combining this solution with Dynamic Time Warping creates a program, that work for both explosive and slow actions.

\subsection{Challenge 4: Translating Similarity into Correctness}
A raw score is informative, but it does not by itself define what counts as a correct technique. A threshold that is appropriate for one action may be too strict or too lenient for another. This problem becomes more severe when different techniques vary in visibility, temporal stability, and sensitivity to joint noise.

The solution implemented in the project is \texttt{calibrate\_label\_thresholds.py}. This script assumes a validation layout of \texttt{datasets/validation/<technique>/correct/*.npy} and \texttt{datasets/validation/<technique>/incorrect/*.npy}. For each technique, it scores validation sequences against the reference bank, searches across candidate thresholds, and chooses the threshold that maximizes F1 score and then balanced accuracy \cite{projectCalibrate}.

\subsection{Challenge 5: Runtime Constraints and Detection Gaps}
The runtime logs show that pose-based reference collection is feasible on CPU hardware, but not trivial. This project has been run using the following specifications:

\begin{itemize}
    \item Processor	AMD Ryzen 9 5900X 12-Core Processor             (3.70 GHz)
    \item Installed RAM	16,0 GB
    \item System type	64-bit operating system, x64-based processor
    \item NVIDIA GeForce GTX 1660 SUPER
    \item NVIDIA Cuda ON
\end{itemize}

In front-kick reference expansion runs, the code processed frames at approximately 10--16 frames per second in many segments, and the logs also contain repeated \textit{no detections} frames \cite{projectFrontKickLog,projectFrontKickStrictLog}. The seed batch run log similarly illustrates the non-negligible cost of continuous pose inference \cite{projectSeedLog}. This is consistent with the literature, which argues that real-time combat analysis is achievable only with careful management of resources and data quality \cite{quinn2024automation,quinn2025yolo}.

The implementation mitigates this through several engineering controls. The scripts expose configurable frame skipping, CPU thread limits, cooldown periods between batch jobs, preflight source checks, headless operation, and a \texttt{fast-mode} preset that disables expensive visual components \cite{projectActionRecognition,projectRunBatch}. Furthermore, CUDA enabled massive parallelism, allowing tasks to run much faster than on traditional CPUs. This has accelerated video processing significantly. 
These measures are not incidental; they are part of the project's answer to practical deployment constraints.


% --- Experimental Setup (moved up) ---
\section{Experimental Setup}
The uploaded Git Repository supports two related evaluation modes. The first is \textit{reference collection and calibration}, where canonical technique examples are captured and validated. The second is \textit{user evaluation}, where an athlete performs a target technique and the system compares that pose sequence against the calibrated reference bank.

Reference examples are stored as NumPy arrays of keypoint sequences in COCO-17 format \cite{projectActionRecognition}. The project documentation specifies a multi-angle folder structure for techniques, while the progress document shows that the intended capture plan consists of 52 total targets. The batch collection script aims by default for four examples per angle, subject to sufficient source diversity \cite{projectRunBatch}. This means the intended experiment design is systematic rather than ad hoc.

For evaluation, \texttt{calibrate\_label\_thresholds.py} requires validation examples for each technique split into \texttt{correct} and \texttt{incorrect} folders \cite{projectCalibrate}. Every validation clip is scored using the same best-reference matching routine used by the trainer, and the threshold is selected automatically from the observed score distribution. This is important because it means the acceptance rule is derived from validation evidence rather than hand tuning.

For an \textit{amateur versus professional} experiment, both groups can perform the same target techniques under comparable recording conditions, after which the system score, threshold decision, and score variance can be compared across groups. For an \textit{expert versus system} experiment, professional practitioner provide an oral score on a 1--100 scale for each attempt, and that score can then be compared against a normalized version of the system's internal score. The code supports the core metrics needed for these experiments.



% --- Results ---
\section{Results}

This section presents the main evaluation results for the implemented system across supported techniques and user groups.

\subsection{Technique Score Summary}

Figure~\ref{fig:score_by_technique} shows the distribution of scores for all supported techniques, aggregated across all evaluated clips. This summary visualization allows for direct comparison of system performance and threshold calibration across techniques, highlighting which actions are more consistently performed above or below the acceptance threshold.


\subsection{Score Distributions by Technique and User}

\begin{figure}[H]
    \centering
    % Replace 'score_distributions_by_technique.png' with your actual plot file
    \includegraphics[width=0.95\textwidth]{score_by_technique_by_technique.png}
    \caption{Score distributions by technique and user group. Each boxplot shows the spread of scores for a given technique, grouped by user and labeled as (Amateur) or (Professional).}
    \label{fig:score_by_technique}
\end{figure}

Figure~\ref{fig:score_by_technique} presents the distribution of system-generated scores for each technique, separated by user and group. The boxplots reveal several key findings:

\begin{itemize}
    \item \textbf{Discriminative Power:} The system distinguishes between amateur and professional users, with professionals generally achieving higher and more consistent scores, especially for techniques like axe kick and front kick.
    \item \textbf{Technique Variability:} Some techniques, such as jab, show a wider spread of scores, indicating greater variability in performance or higher sensitivity in the scoring algorithm.
    \item \textbf{Outliers and Consistency:} The presence of outliers suggests occasional detection errors or natural performance variation. Professionals tend to have fewer outliers and narrower score ranges.
\end{itemize}

These results demonstrate that the program’s scoring system is sensitive to both user skill and technique, aligning with the intended design of providing meaningful, technique-specific feedback.


% --- Analysis ---
\section{Analysis}
At a qualitative level, the reported examples are consistent with the intended trainer logic. Windows above threshold are marked as correct, whereas lower-scoring windows trigger corrective feedback and, where configured, ghost-pose guidance. The figure-based jab comparison also serves as a screenshot-based demonstration of the system output, showing the tracked skeleton, bounding box, activity label, and score overlay while visually highlighting differences in extension and guard position across examples. These observations are compatible with the system design, but they remain clip-level observations rather than participant-level findings.

\subsection*{Interpretation of Technique Differences and Score Distributions}

The score distribution plots highlight important differences between techniques. Techniques such as jab exhibit a broader range of scores, which may reflect greater difficulty, more variation in user execution, or increased sensitivity in the scoring logic. In contrast, techniques like axe kick and front kick show professionals consistently outperforming amateurs, suggesting that the system effectively captures skill differences for these actions.

It is important to note that most results achieve relatively high scores. This is partly due to the use of a generic fighting stance as the reference, which is a standard pose for executing most techniques. As a result, the fighting stance may bias the scores upward, making it easier for users to achieve high marks when starting from this position.

However, when the program is used to interpret techniques performed outside of a fighting stance, the score and feedback indicators become more reliable and discriminative. Therefore, the mean score in these plots does not fully capture the quality of the technique itself, but is also influenced by the initial stance and reference selection. This should be considered when interpreting the results and comparing across techniques.

Overall, the plots confirm that the system can differentiate between user groups and techniques, but also highlight the need for careful reference selection and further calibration to ensure that scores reflect true technique quality rather than just conformity to a standard pose.

% --- Discussion ---
\section{Discussion}
At a qualitative level, the reported examples are consistent with the intended trainer logic. Windows above threshold are marked as correct, whereas lower-scoring windows trigger corrective feedback and, where configured, ghost-pose guidance. The figure-based jab comparison also serves as a screenshot-based demonstration of the system output, showing the tracked skeleton, bounding box, activity label, and score overlay while visually highlighting differences in extension and guard position across examples. These observations are compatible with the system design, but they remain clip-level observations rather than participant-level findings.

\subsection*{Interpretation of Technique Differences and Score Distributions}

The score distribution plots highlight important differences between techniques. Techniques such as jab exhibit a broader range of scores, which may reflect greater difficulty, more variation in user execution, or increased sensitivity in the scoring logic. In contrast, techniques like axe kick and front kick show professionals consistently outperforming amateurs, suggesting that the system effectively captures skill differences for these actions.

It is important to note that most results achieve relatively high scores. This is partly due to the use of a generic fighting stance as the reference, which is a standard pose for executing most techniques. As a result, the fighting stance may bias the scores upward, making it easier for users to achieve high marks when starting from this position.

However, when the program is used to interpret techniques performed outside of a fighting stance, the score and feedback indicators become more reliable and discriminative. Therefore, the mean score in these plots does not fully capture the quality of the technique itself, but is also influenced by the initial stance and reference selection. This should be considered when interpreting the results and comparing across techniques.

Overall, the plots confirm that the system can differentiate between user groups and techniques, but also highlight the need for careful reference selection and further calibration to ensure that scores reflect true technique quality rather than just conformity to a standard pose.

\begin{figure}[h!]
\centering
% Top row
\begin{minipage}{0.48\textwidth}
    \centering
    \includegraphics[width=\linewidth, height=6cm, keepaspectratio]{Jab_FP_pro.png}
\end{minipage}
\hfill
\begin{minipage}{0.48\textwidth}
    \centering
    \includegraphics[width=\linewidth, height=6cm, keepaspectratio]{Jab_FP_ext_Pro.png}
\end{minipage}

\vspace{0.5cm}
% Bottom row
\begin{minipage}{0.48\textwidth}
    \centering
    \includegraphics[width=\linewidth, height=6cm, keepaspectratio]{Jab_Ama_Guard_Good.png}
\end{minipage}
\hfill
\begin{minipage}{0.48\textwidth}
    \centering
    \includegraphics[width=\linewidth, height=6cm, keepaspectratio]{Jab_Ama_Guard.png}
\end{minipage}

\caption{Illustrative jab comparison. The examples visually highlight differences in extension and guard position, which correspond to the kinds of deviations the trainer is designed to score and comment on.}
\label{fig:jab_comparison}
\end{figure}

Figure~\ref{fig:roundhouse_correction} presents a qualitative example of corrective feedback for a roundhouse kick. In this sequence, the system indicates that the kicking knee is driven too directly toward the target instead of first moving outward and upward during the development phase. Although this straighter path may feel more natural to an amateur performer, it is technically inefficient because it reduces rotational force generation and contributes to loss of balance. A visible consequence in this example is that the left guard drops during execution, suggesting that defensive posture is sacrificed to stabilize the movement. The final frame includes a highlighted reference overlay in order to make the intended correction easier to interpret.

\begin{figure}[htbp]
    \centering
    
\begin{subfigure}[t]{0.32\textwidth}
    \centering
    \includegraphics[width=\linewidth,height=6.2cm]{Dicte_Roundhouse1.png}
    \caption{Preparation}
\end{subfigure}
\hfill
\begin{subfigure}[t]{0.32\textwidth}
    \centering
    \includegraphics[width=\linewidth,height=6.2cm]{Dicte_Roundhouse3.png}
    \caption{Development}
\end{subfigure}
\hfill
\begin{subfigure}[t]{0.32\textwidth}
    \centering
    \includegraphics[width=\linewidth,height=6.2cm]{Dicte_Roundhouse6.png}
    \caption{Execution}
\end{subfigure}
    
    \caption{Qualitative example of corrective feedback for a roundhouse kick. The sequence shows that the kicking knee follows a path that is too direct, moving inward and forward rather than first lifting outward and upward during the chamber phase. This is a technically important error, since it reduces rotational loading, limits power generation, slows the kick, and creates visible imbalance. In this example, the imbalance is accompanied by lowering of the left guard, indicating compensation in the upper body. In the final frame, the reference overlay has been manually highlighted to make the intended correction easier to distinguish. The example illustrates that the system can reveal meaningful technical deficiencies in movement trajectory and posture, not only in the final extended position.}
    \label{fig:roundhouse_correction}
\end{figure}

\subsection*{Program Score Across Evaluated Runs}
A clip-level comparison was conducted across the evaluated runs using the latest structured run metrics. In this comparison, the program score is derived from the automated scoring pipeline, while the acceptance statistics provide a complementary view of how consistently temporal windows within each clip were judged as technically acceptable. The purpose of this comparison is not to claim completed external validation, but to examine how the system behaves across multiple techniques and performers.

Tables~\ref{tab:albert_run_summary}--\ref{tab:jimmy_run_summary} summarize the clip-level results grouped by performer. Here, \textit{Windows} denotes the number of evaluated temporal segments extracted from a clip, \textit{Accepted} denotes the number of segments classified as correct by the trainer, and \textit{Acceptance rate} expresses the proportion of accepted windows. The \textit{Program Score} is the aggregate score produced by the system, while the standard deviation indicates the spread of window-level scores within each clip.

\begin{table}[H]
\centering
\caption{Run-level summary for Albert.}
\label{tab:albert_run_summary}
\resizebox{\textwidth}{!}{
\begin{tabular}{lccccc}
\toprule
Technique & Windows & Accepted & Acceptance rate (\%) & Program Score & Std.\ dev. \\
\midrule
Axe Kick   & 582 & 170 & 29.2 & 64.75 & 9.65 \\
Cross      & 582 & 140 & 24.1 & 58.18 & 12.58 \\
Elbow      & 172 & 73  & 42.4 & 66.15 & 12.48 \\
Front Kick & 266 & 144 & 54.1 & 68.02 & 11.39 \\
Jab        & 315 & 196 & 62.2 & 69.63 & 9.64 \\
Uppercut   & 417 & 113 & 27.1 & 60.83 & 11.16 \\
\bottomrule
\end{tabular}
}
\end{table}

\begin{table}[H]
\centering
\caption{Run-level summary for Dicte.}
\label{tab:dicte_run_summary}
\resizebox{\textwidth}{!}{
\begin{tabular}{lccccc}
\toprule
Technique & Windows & Accepted & Acceptance rate (\%) & Program Score & Std.\ dev. \\
\midrule
Axe Kick        & 1046 & 295 & 28.2 & 62.33 & 10.45 \\
Cross           & 780  & 142 & 18.2 & 56.40 & 10.83 \\
Elbow           & 801  & 234 & 29.2 & 61.76 & 10.48 \\
Front Kick      & 771  & 56  & 7.3  & 53.51 & 10.19 \\
Hook            & 826  & 226 & 27.4 & 60.28 & 11.05 \\
Jab             & 638  & 197 & 30.9 & 60.48 & 11.01 \\
Knee Kick       & 869  & 307 & 35.3 & 63.79 & 13.48 \\
Roundhouse Kick & 1104 & 228 & 20.7 & 60.69 & 10.14 \\
Uppercut        & 762  & 98  & 12.9 & 58.54 & 9.16 \\
\bottomrule
\end{tabular}
}
\end{table}

\begin{table}[H]
\centering
\caption{Run-level summary for Jimmy.}
\label{tab:jimmy_run_summary}
\resizebox{\textwidth}{!}{
\begin{tabular}{lccccc}
\toprule
Technique & Windows & Accepted & Acceptance rate (\%) & Program Score & Std.\ dev. \\
\midrule
Axe Kick        & 442  & 440 & 99.5 & 89.17 & 4.14 \\
Cross           & 804  & 125 & 15.5 & 56.86 & 10.87 \\
Elbow           & 626  & 189 & 30.2 & 66.11 & 10.92 \\
Front Kick      & 1540 & 499 & 32.4 & 64.91 & 14.16 \\
Hook            & 548  & 336 & 61.3 & 71.51 & 15.20 \\
Jab             & 302  & 277 & 91.7 & 80.73 & 7.98 \\
Knee Kick       & 700  & 424 & 60.6 & 75.11 & 17.39 \\
Roundhouse Kick & 885  & 263 & 29.7 & 61.81 & 18.96 \\
Uppercut        & 346  & 228 & 65.9 & 74.63 & 12.26 \\
\bottomrule
\end{tabular}
}
\end{table}

% --- Conclusion ---
\section{Conclusion}

This thesis presented a computer-vision-based system for automated combat-sports technique analysis. The implemented approach uses pose estimation, temporal sequence comparison, multi-angle reference matching, structured run metrics, and threshold-based scoring to evaluate the quality of a performed action. Unlike a conventional action-recognition system, the goal was not only to identify a movement, but to assess how well it was executed and to provide feedback that is meaningful in a training context.

The results show that this objective is achievable at a prototype level. The system can extract pose sequences from video, compare them against stored references, generate numerical scores, classify windows as accepted or rejected, and provide corrective cues linked to observable technical issues. In the jab examples reported in this thesis, the system produced score patterns that broadly reflected differences in execution quality, and its output was directionally consistent with human judgement in the limited oral-score comparison.

The project therefore makes two main contributions. First, it demonstrates a complete technical pipeline for reference-based martial-arts assessment. Second, it shows that pose-driven scoring can move beyond pure action classification toward qualitative evaluation of technique. These contributions are meaningful because they establish both a working software artefact and a methodological basis for later empirical validation.

At the same time, the thesis also makes clear that the current work should be interpreted as an initial validation rather than a finished assessment framework. The present evidence is limited in scale, threshold calibration is not yet complete for the available snapshot, and expert comparison has only been demonstrated on a small illustrative set. As a result, the strongest defensible conclusion is not that the system already matches expert judgement reliably, but that it provides a credible and technically coherent foundation for doing so in future work.

% --- Future Work ---
\section{Future Work}
Several extensions would strengthen the system substantially. The first priority is to expand and clean the reference library so that each target technique is represented by multiple high-quality examples from controlled viewpoints. The second is to build labeled validation sets for each technique, allowing the implemented calibration pipeline to produce defensible thresholds and report standard metrics such as precision, recall, F1 score, and balanced accuracy. The third is to conduct a larger expert-evaluation study in which multiple coaches score the same clips and the system outputs are compared against those ratings using formal agreement measures.

A further direction for future work is to improve robustness in real-world footage. This includes handling more challenging viewpoints, reducing the effect of missing detections, and refining the temporal segmentation of short actions. It would also be useful to extend the evaluation beyond jab techniques to a broader set of punches and kicks, since that would test whether the scoring framework generalizes across movements with different biomechanics.

In conclusion, the thesis shows that automated combat-sports technique assessment using computer vision is feasible and practically promising. While stronger empirical validation is still required, the implemented system already demonstrates the essential components of a pose-based digital trainer and provides a strong basis for continued development and evaluation.

\printbibliography

\end{document}
